\subsection{Scalar consequences of the common inward budget}
\label{ext:inward-scalar}
Some finite count domains use a stronger absolute-third-moment floor or
noise ceiling before their five-dimensional subdivision. The following
calculation supplies those scalar restrictions. It uses the same sign
conditioning as Section~\ref{ext:count-filters}.

Reflect only for the scalar calculation, so that
$0<p\le1/2$, $q=1-p$, $v=pq$. This preserves $\rho$, total
conditional variance $s$ and total inward displacement. The count
probability, its phase and its original orientation are not reflected.
Suppose the inward endpoints obey $a_0+a_1\le d$ with $0<d\le p$.
The conditional third-moment inequalities and Cauchy's inequality give
\begin{equation}\label{ext:inward-I}
 \rho\ge I(p,s,d):=
 \frac{v(1-2v)+(3p-d)s+s^2/(qd)}{(v+s)^{3/2}}.
\end{equation}
Indeed, writing $b=qs_0$, $c=ps_1$ and $b+c=s$, the absolute-cubic
numerator is at least
\[
 v(1-2v)+3pb+3qc-a_0b-a_1c
                       +\frac{b^2}{qa_0}+\frac{c^2}{pa_1}.
\]
Its linear part is at least $(3p-d)s$, while its quadratic part is
at least $(b+c)^2/(q(a_0+a_1))\ge s^2/(qd)$.
Vanishing inward endpoints force the corresponding centered variance
to vanish; those terms are handled before division. If $d=0$, then
$s=0$ directly.

The derivative has the exact numerator
\begin{equation}\label{ext:inward-I-derivative}
 (v+s)^{5/2}I_s
 =v(6p-3p^2-d-3/2)
   +\left\{\frac{2p}{d}-\frac{3p-d}{2}\right\}s
   +\frac{s^2}{2qd}.
\end{equation}
Consequently, on $p\in[l,u]\subset(0,1/2]$, the two gates
$d\le l$ and $f:=6l-3l^2-d-3/2\ge0$ imply
\[
 \rho\ge\rho_B(p):=\frac{1-2p(1-p)}{\sqrt{p(1-p)}}
                  \ge\rho_B(u).
\]
The last function is decreasing on this half-interval. If a previously
acquired ceiling is $\rho\le R$, a point with
\[
 R^2p(1-p)-(1-2p(1-p))^2<0
\]
is impossible. This polynomial increases on $0<p\le1/2$, because
its derivative is $(1-2p)(R^2+4-8p(1-p))\ge0$.
Exact bisection may therefore raise the closed lower probability
endpoint to a known infeasible point; retaining that endpoint is a
conservative enlargement.

When $f>0$ and $s\le S$, formula~\eqref{ext:inward-I-derivative}
also gives
\[
 I_s\ge
 \frac{l(1-l)f}{[u(1-u)+S]^{5/2}}.
\]
Integration from zero supplies a new variance ceiling. Directed square
roots and outward division make every finite update valid separately.

A useful weaker floor survives failure of the monotonicity gate.
Assume still $0<d\le l$ and set
\[
 a=(3/2-6l+3l^2+d)_+,
 \qquad B=2l/d-3u/2\ge5/4.
\]
The only possible decrease of $I$ is bounded below by
$-(va-Bs)_+/v^{5/2}$. Integrating this negative triangle gives
\begin{equation}\label{ext:inward-deficit}
 \rho\ge\rho_B(u)-\frac{a^2}{2B\sqrt{l(1-l)}}.
\end{equation}
In particular, failure of the monotonicity gate must not be replaced
by the unsupported assertion $\rho\ge\rho_B(p)$.

There is a second variance contraction which works with either sign
of $f$. The global Taylor upper bound for $(v+s)^{3/2}$, combined
with \eqref{ext:inward-I}, gives the necessary inequality
\begin{equation}\label{ext:scalar-noise-quadratic}
 Gs^2+fs\le\epsilon\,[u(1-u)+S]^{3/2},
\end{equation}
where
\[
 \epsilon=(R-\rho_B(u))_+,\qquad
 G=\frac1{(1-l)d}-\frac{3(1-2l(1-l))}{8l(1-l)}
       \ge\frac5{8l(1-l)}+\frac34>0.
\]
To see this, subtract $\rho_B(p)(v+s)^{3/2}$ from the lower
numerator in \eqref{ext:inward-I}, use the quadratic Taylor upper
bound on the subtracted term, and then use $p\ge l$ and
$\rho_B(p)\ge\rho_B(u)$. The positive part in $\epsilon$ is
necessary when $f<0$. The larger real root of
\eqref{ext:scalar-noise-quadratic}, enclosed from above, is a valid
new upper bound for $s$, even for negative $f$.

Finally suppose the original joint resource has the stronger form
\[
 d_{\rm actual}\le d_0(1-s/S_0),\qquad S_0=\Lambda/n.
\]
To test a proposed upper endpoint $a\in(0,S_0)$, assume $s\ge a$.
Then $d_{\rm actual}\le d_0(1-a/S_0)$, so the preceding quadratic
calculation can be repeated with this smaller inward bound. If its
resulting upper bound is strictly less than $a$, the whole region
$s\ge a$ is impossible. Otherwise the probe proves no lower bound
on $s$. A finite sequence of such hypothetical tests is sound
without any convergence assertion. Throughout it, the original
$S_0$ remains in the joint budget; a separately improved noise
ceiling clips that budget and does not replace its intercept.
